\subsection{Transformations of graphs}

A known graph can generate a family of new graphs without starting from zero. If the graph of \(y=f(x)\) is known, then
\[
y=f(x)+k
\]
shifts it vertically by \(k\), while
\[
y=f(x-h)
\]
shifts it horizontally by \(h\). Multiplication changes scale:
\[
y=af(x)
\]
stretches or compresses vertical distances, and a negative value of \(a\) also reflects the graph across the horizontal axis.

\begin{examplebox}[Building graphs from \(x^2\)]
Starting from \(f(x)=x^2\), the function
\[
g(x)=-(x-2)^2+3
\]
is obtained by shifting the parabola two units to the right, reflecting it across the horizontal axis, and shifting it three units upward. Its vertex is therefore \((2,3)\), and it opens downward.
\end{examplebox}

\begin{interpretationbox}[Why transformations matter]
The purpose is not to memorize isolated pictures. Transformations let us read parameters in a formula as geometric instructions. This becomes useful later when parameters change the position of zeros, extrema, or asymptotes.
\end{interpretationbox}

\begin{summarybox}[Graph transformations]
Separate horizontal changes inside the argument from vertical changes outside the function. State each transformation and then describe its visible effect on the graph.
\end{summarybox}
